% Options for packages loaded elsewhere
% Options for packages loaded elsewhere
\PassOptionsToPackage{unicode}{hyperref}
\PassOptionsToPackage{hyphens}{url}
\PassOptionsToPackage{dvipsnames,svgnames,x11names}{xcolor}
%
\documentclass[
  spanish,
  10pt,
  letterpaper,
  twocolumn]{article}
\usepackage{xcolor}
\usepackage[top=1.0cm,bottom=1.5cm,left=1.5cm,right=1.5cm]{geometry}
\usepackage{amsmath,amssymb}
\setcounter{secnumdepth}{-\maxdimen} % remove section numbering
\usepackage{iftex}
\ifPDFTeX
  \usepackage[T1]{fontenc}
  \usepackage[utf8]{inputenc}
  \usepackage{textcomp} % provide euro and other symbols
\else % if luatex or xetex
  \usepackage{unicode-math} % this also loads fontspec
  \defaultfontfeatures{Scale=MatchLowercase}
  \defaultfontfeatures[\rmfamily]{Ligatures=TeX,Scale=1}
\fi
\usepackage{lmodern}
\ifPDFTeX\else
  % xetex/luatex font selection
  \setmainfont[]{Fira Sans}
\fi
% Use upquote if available, for straight quotes in verbatim environments
\IfFileExists{upquote.sty}{\usepackage{upquote}}{}
\IfFileExists{microtype.sty}{% use microtype if available
  \usepackage[]{microtype}
  \UseMicrotypeSet[protrusion]{basicmath} % disable protrusion for tt fonts
}{}
\usepackage{setspace}
\makeatletter
\@ifundefined{KOMAClassName}{% if non-KOMA class
  \IfFileExists{parskip.sty}{%
    \usepackage{parskip}
  }{% else
    \setlength{\parindent}{0pt}
    \setlength{\parskip}{6pt plus 2pt minus 1pt}}
}{% if KOMA class
  \KOMAoptions{parskip=half}}
\makeatother
% Make \paragraph and \subparagraph free-standing
\makeatletter
\ifx\paragraph\undefined\else
  \let\oldparagraph\paragraph
  \renewcommand{\paragraph}{
    \@ifstar
      \xxxParagraphStar
      \xxxParagraphNoStar
  }
  \newcommand{\xxxParagraphStar}[1]{\oldparagraph*{#1}\mbox{}}
  \newcommand{\xxxParagraphNoStar}[1]{\oldparagraph{#1}\mbox{}}
\fi
\ifx\subparagraph\undefined\else
  \let\oldsubparagraph\subparagraph
  \renewcommand{\subparagraph}{
    \@ifstar
      \xxxSubParagraphStar
      \xxxSubParagraphNoStar
  }
  \newcommand{\xxxSubParagraphStar}[1]{\oldsubparagraph*{#1}\mbox{}}
  \newcommand{\xxxSubParagraphNoStar}[1]{\oldsubparagraph{#1}\mbox{}}
\fi
\makeatother


\usepackage{longtable,booktabs,array}
\newcounter{none} % for unnumbered tables
\usepackage{calc} % for calculating minipage widths
% Correct order of tables after \paragraph or \subparagraph
\usepackage{etoolbox}
\makeatletter
\patchcmd\longtable{\par}{\if@noskipsec\mbox{}\fi\par}{}{}
\makeatother
% Allow footnotes in longtable head/foot
\IfFileExists{footnotehyper.sty}{\usepackage{footnotehyper}}{\usepackage{footnote}}
\makesavenoteenv{longtable}
\usepackage{graphicx}
\makeatletter
\newsavebox\pandoc@box
\newcommand*\pandocbounded[1]{% scales image to fit in text height/width
  \sbox\pandoc@box{#1}%
  \Gscale@div\@tempa{\textheight}{\dimexpr\ht\pandoc@box+\dp\pandoc@box\relax}%
  \Gscale@div\@tempb{\linewidth}{\wd\pandoc@box}%
  \ifdim\@tempb\p@<\@tempa\p@\let\@tempa\@tempb\fi% select the smaller of both
  \ifdim\@tempa\p@<\p@\scalebox{\@tempa}{\usebox\pandoc@box}%
  \else\usebox{\pandoc@box}%
  \fi%
}
% Set default figure placement to htbp
\def\fps@figure{htbp}
\makeatother



\ifLuaTeX
\usepackage[bidi=basic,shorthands=off]{babel}
\else
\usepackage[bidi=default,shorthands=off]{babel}
\fi
\ifPDFTeX
\else
\babelfont{rm}[]{Fira Sans}
\fi
\ifLuaTeX
  \usepackage{selnolig} % disable illegal ligatures
\fi


\setlength{\emergencystretch}{3em} % prevent overfull lines

\providecommand{\tightlist}{%
  \setlength{\itemsep}{0pt}\setlength{\parskip}{0pt}}



 


\usepackage{tikz}
\makeatletter
\@ifpackageloaded{caption}{}{\usepackage{caption}}
\AtBeginDocument{%
\ifdefined\contentsname
  \renewcommand*\contentsname{Tabla de contenidos}
\else
  \newcommand\contentsname{Tabla de contenidos}
\fi
\ifdefined\listfigurename
  \renewcommand*\listfigurename{Listado de Figuras}
\else
  \newcommand\listfigurename{Listado de Figuras}
\fi
\ifdefined\listtablename
  \renewcommand*\listtablename{Listado de Tablas}
\else
  \newcommand\listtablename{Listado de Tablas}
\fi
\ifdefined\figurename
  \renewcommand*\figurename{Figura}
\else
  \newcommand\figurename{Figura}
\fi
\ifdefined\tablename
  \renewcommand*\tablename{Tabla}
\else
  \newcommand\tablename{Tabla}
\fi
}
\@ifpackageloaded{float}{}{\usepackage{float}}
\floatstyle{ruled}
\@ifundefined{c@chapter}{\newfloat{codelisting}{h}{lop}}{\newfloat{codelisting}{h}{lop}[chapter]}
\floatname{codelisting}{Listado}
\newcommand*\listoflistings{\listof{codelisting}{Listado de Listados}}
\makeatother
\makeatletter
\makeatother
\makeatletter
\@ifpackageloaded{caption}{}{\usepackage{caption}}
\@ifpackageloaded{subcaption}{}{\usepackage{subcaption}}
\makeatother
\usepackage{bookmark}
\IfFileExists{xurl.sty}{\usepackage{xurl}}{} % add URL line breaks if available
\urlstyle{same}
% fallback for those not using the hyperref driver hyperxmp:
\makeatletter
\@ifundefined{xmpquote}{\newcommand{\xmpquote}[1]{#1}}{}
\makeatother
\hypersetup{
  pdftitle={Práctica No.~2 -- Relaciones Hídricas},
  pdflang={es},
  colorlinks=true,
  linkcolor={blue},
  filecolor={Maroon},
  citecolor={Blue},
  urlcolor={Blue},
  pdfcreator={LaTeX via pandoc}}


\title{\vspace{-1.5em}{\Huge \textbf{Práctica No.~2 -- Relaciones
Hídricas}}\vspace{-1.5em}}
\date{}
\begin{document}
\maketitle


\setstretch{1.2}
\textbf{Fecha}: \hrulefill \hfill \textbf{Grupo}: \hrulefill

\textbf{Nombre:} \hrulefill

\textbf{Objetivo:} Reconocer la manifestación de difusión, ósmosis y
plasmólisis en células vegetales; observar cambios de turgencia en
medios hipertónicos, hipotónicos e isotónicos; y determinar el potencial
hídrico de un tejido vegetal mediante los métodos gravimétrico y
densimétrico.

\subsection{1. Materiales}\label{materiales}

\begin{itemize}
\tightlist
\item
  Hojas de \emph{Elodea} (\emph{Egeria densa})
\item
  Tubérculo de papa (\emph{Solanum tuberosum})
\item
  Kit de disección
\item
  Portaobjetos y cubreobjetos
\item
  Regla
\end{itemize}

\subsection{2. Aspectos teóricos sobre relaciones
hídricas}\label{aspectos-teuxf3ricos-sobre-relaciones-huxeddricas}

Las relaciones hídricas en las plantas se refieren al movimiento y
distribución del agua dentro del organismo vegetal y su interacción con
el ambiente circundante. El potencial hídrico (\(\Psi\)) es una medida
que determina la dirección y la velocidad del flujo de agua en las
células vegetales, estando determinado fundamentalmente por el potencial
osmótico (\(\Psi_s\)) y el potencial de presión (\(\Psi_p\)).

Defina brevemente los siguientes conceptos:

\subsubsection{\texorpdfstring{1. Potencial hídrico (\(\Psi\)) y sus
componentes (\(\Psi_s\),
\(\Psi_p\)):}{1. Potencial hídrico (\textbackslash Psi) y sus componentes (\textbackslash Psi\_s, \textbackslash Psi\_p):}}\label{potencial-huxeddrico-psi-y-sus-componentes-psi_s-psi_p}

\hrulefill

\hrulefill

\hrulefill

\hrulefill

\hrulefill

\subsubsection{2. Plasmólisis y plasmólisis
incipiente:}\label{plasmuxf3lisis-y-plasmuxf3lisis-incipiente}

\hrulefill

\hrulefill

\hrulefill

\hrulefill

\hrulefill

\subsubsection{\texorpdfstring{3. Ecuación de Van't Hoff
(\(\Psi = -RTC\)) y explicación de sus
variables:}{3. Ecuación de Van't Hoff (\textbackslash Psi = -RTC) y explicación de sus variables:}}\label{ecuaciuxf3n-de-vant-hoff-psi--rtc-y-explicaciuxf3n-de-sus-variables}

\hrulefill

\hrulefill

\hrulefill

\hrulefill

\hrulefill

\subsection{3. Cambios de turgencia
celular}\label{cambios-de-turgencia-celular}

\begin{enumerate}
\def\labelenumi{\arabic{enumi}.}
\tightlist
\item
  Tome varias hojas de \emph{Elodea} y colóquelas en cajas de Petri con
  soluciones de sacarosa de 0.05, 0.15, 0.25, 0.50 y 0.75 M durante 10
  minutos.
\item
  Realice un montaje húmedo con portaobjetos y cubreobjetos (comenzando
  por observar las muestras de 0.05 M y 0.75 M).
\item
  Observe en el microscopio óptico 25 células por tratamiento y anote
  cuántas aparecen turgentes y cuántas plasmolizadas.
\end{enumerate}

\begin{center}
\begin{tikzpicture}
  \draw[thick] (0,0) circle (2.8cm);
  \node[below] at (0,-2.9) {\textbf{Muestra con 0.05 M}};
\end{tikzpicture}
\end{center}

\begin{center}
\begin{tikzpicture}
  \draw[thick] (0,0) circle (2.8cm);
  \node[below] at (0,-2.9) {\textbf{Muestra con 0.75 M}};
\end{tikzpicture}
\end{center}

Describa los cambios celulares observados en la membrana plasmática
respecto a la pared celular en cada concentración:

\hrulefill

\hrulefill

\hrulefill

\hrulefill

\hrulefill

\hrulefill

¿En cuál solución de sacarosa el 50\% de las células presenta
plasmólisis incipiente (concentración isotónica)?

\hrulefill

\hrulefill

\hrulefill

\hrulefill

\subsection{4. Potencial hídrico en tejidos
vegetales}\label{potencial-huxeddrico-en-tejidos-vegetales}

\textbf{Temperatura ambiental del laboratorio:}
\hrulefill \(^\circ\text{C}\) (\(\text{K} =\) \hrulefill)

En un vaso de precipitado con agua introduzca un termómetro y registre
la temperatura estabilizada.

\subsubsection{a. Método gravimétrico (cambio de
masa)}\label{a.-muxe9todo-gravimuxe9trico-cambio-de-masa}

\begin{enumerate}
\def\labelenumi{\arabic{enumi}.}
\item
  En una batería de 5 tubos de ensayo, adicione 10 mL de cada solución
  de sacarosa (0.05 M a 0.75 M).
\item
  Corte 5 cilindros de papa (\emph{Solanum tuberosum}) de
  aproximadamente 5 cm de largo.
\item
  Pese cada cilindro, registre la masa inicial, introdúzcalos en cada
  tubo durante 1 hora, seque el exceso de humedad y peselos nuevamente.
\end{enumerate}

\begin{tabular}{c|c|p{1.5cm}|p{1.5cm}|p{1.5cm}}
\hline
Tubo & Conc. (M) & Masa Inicial (g) & Masa Final (g) & Cambio Masa (\%) \\
\hline
1 & 0.05 &  &  &  \\
\hline
2 & 0.15 &  &  &  \\
\hline
3 & 0.25 &  &  &  \\
\hline
4 & 0.50 &  &  &  \\
\hline
5 & 0.75 &  &  &  \\
\hline
\end{tabular}

Calcule el potencial hídrico para la solución en la que no hubo cambio
de masa expresado en MPa (\(\Psi = -RTC\)):

\hrulefill

\hrulefill

Interprete los resultados obtenidos mediante el método gravimétrico:

\hrulefill

\hrulefill

\hrulefill

\hrulefill

\hrulefill

\subsubsection{b. Método densimétrico
(Chardakov)}\label{b.-muxe9todo-densimuxe9trico-chardakov}

\begin{enumerate}
\def\labelenumi{\arabic{enumi}.}
\item
  Agregue una gota de azul de metileno a las soluciones donde reposaron
  las papas y agite con cuidado.
\item
  Extraiga una gota con pipeta y trasfiérala a los tubos de la segunda
  batería testigo (solución sin colorante). Evite agitar o intentar
  disolver la gota en la solución sin colorante.
\item
  Observe si la gota desciende al fondo, sube hacia la superficie o
  permanece estacionaria.
\end{enumerate}

\begin{center}
\begin{tikzpicture}[scale=0.85]
  \foreach \x/\ydrop/\lbl in {
    0/2.1/{\textbf{La gota sube:} la solución se diluyó, puesto que el tejido vegetal pierde agua.},
    2.7/1.3/{\textbf{La gota permanece o se difunde:} los potenciales en el tejido y la solución son iguales.},
    5.4/0.55/{\textbf{La gota baja:} la solución se concentró y aumentó su densidad, transfirió agua al tejido vegetal.}
  } {
    % Borde superior acampanado y cuerpo del tubo de ensayo
    \draw[thick] (\x-0.08, 3.2) -- (\x, 3.0) -- (\x, 0.5) 
                 arc (180:360:0.45) -- (\x+0.9, 3.0) -- (\x+0.98, 3.2);

    % Menisco de la solución
    \draw[thick] (\x, 2.5) .. controls (\x+0.25, 2.2) and (\x+0.65, 2.2) .. (\x+0.9, 2.5);

    % Gota colorante (forma de lágrima)
    \fill[blue!85!black] (\x+0.45, \ydrop+0.18) 
      .. controls (\x+0.22, \ydrop+0.02) and (\x+0.22, \ydrop-0.25) .. (\x+0.45, \ydrop-0.25) 
      .. controls (\x+0.68, \ydrop-0.25) and (\x+0.68, \ydrop+0.02) .. cycle;

    % Caja descriptiva inferior (Ajustada para evitar colisiones)
    \node[draw=black, fill=gray!12, rounded corners=4pt, inner sep=3pt, text width=2.0cm, align=left, font=\fontsize{6}{7}\selectfont, anchor=north] 
      at (\x+0.45, -0.2) {\lbl};
  }
\end{tikzpicture}
\end{center}

\begin{tabular}{c|c|p{6em}|p{7em}}
\hline
Tubo & Conc. (M) & Movimiento Gota & Interpretación \\
\hline
1 & 0.05 &  &  \\
\hline
2 & 0.15 &  &  \\
\hline
3 & 0.25 &  &  \\
\hline
4 & 0.50 &  &  \\
\hline
5 & 0.75 &  &  \\
\hline
\end{tabular}

Interprete los resultados obtenidos mediante el método densimétrico:

\hrulefill

\hrulefill

\hrulefill

\hrulefill

\hrulefill

\subsection{5. Actividades
complementarias}\label{actividades-complementarias}

\begin{enumerate}
\def\labelenumi{\arabic{enumi}.}
\tightlist
\item
  Explique el concepto de presión osmótica y su importancia biológica.
\end{enumerate}

\hrulefill

\hrulefill

\hrulefill

\hrulefill

\hrulefill

\hrulefill

\hrulefill

\hrulefill

\hrulefill

\hrulefill

\begin{enumerate}
\def\labelenumi{\arabic{enumi}.}
\setcounter{enumi}{1}
\tightlist
\item
  ¿Cómo se modifica el potencial hídrico celular y cada uno de sus
  componentes cuando las células están en contacto con medios iso, hiper
  e hipotónicos?
\end{enumerate}

\hrulefill

\hrulefill

\hrulefill

\hrulefill

\hrulefill

\hrulefill

\hrulefill

\hrulefill

\hrulefill

\hrulefill

\begin{enumerate}
\def\labelenumi{\arabic{enumi}.}
\setcounter{enumi}{2}
\tightlist
\item
  ¿Qué ventajas y desventajas tiene utilizar el método densimétrico para
  medir el potencial hídrico?
\end{enumerate}

\hrulefill

\hrulefill

\hrulefill

\hrulefill

\hrulefill

\hrulefill

\hrulefill

\hrulefill

\hrulefill

\begin{enumerate}
\def\labelenumi{\arabic{enumi}.}
\setcounter{enumi}{3}
\tightlist
\item
  ¿Cuáles son los efectos del potencial hídrico en los procesos
  fisiológicos de las plantas? Explique.
\end{enumerate}

\hrulefill

\hrulefill

\hrulefill

\hrulefill

\hrulefill

\hrulefill

\hrulefill

\hrulefill

\hrulefill

\hrulefill

\subsection{6. Conclusión}\label{conclusiuxf3n}

\hrulefill

\hrulefill

\hrulefill

\hrulefill

\hrulefill

\hrulefill

\hrulefill

\subsection{7. Bibliografía}\label{bibliografuxeda}

\begin{itemize}
\item
  Azcón-Bieto, J., \& Talón, M. (2008). \emph{Fundamentos de Fisiología
  Vegetal}. McGraw-Hill Interamericana.
\item
  Casselles O., A. (2013). \emph{Biología práctica: Guías de trabajo}.
  Colombia.
\item
  Melgarejo, L. M. (Ed.). (2010). \emph{Experimentos en fisiología
  vegetal}. Universidad Nacional de Colombia.
\item
  Taiz, L., \& Zeiger, E. (2002). \emph{Plant Physiology} (3rd ed.).
  Sinauer Associates Ltd.
\end{itemize}




\end{document}
